\section{Experimental Setup and Edge Metrics}
\label{sec:experimental-setup}

All methodology described in this section -- sample sizes, hardware, and measurement procedures --
refers to offline, bench-level evaluation run against synthetically generated sensor data, static
labeled images, and scripted or synthetic scenarios executed against the codebase. None of it refers
to live sensor telemetry from an operating recirculating aquaculture facility; no such facility or
live sensor feed was used anywhere in this paper, and no physical water-quality sensor hardware
exists for this system at all -- every dissolved-oxygen, temperature, and other water-chemistry
value referenced anywhere in this paper, including in Section~\ref{sec:results}'s fabrication
analysis, was synthetically generated for a scripted test scenario, not recorded from a real sensor.
This is a deliberate scope choice, not an incidental one: absent physical sensor hardware, the
purpose of this evaluation is to demonstrate the architectural feasibility of running this system's
multi-component pipeline -- LLM, vision, retrieval, and rule engine together -- and specifically to
test whether its space and memory footprint is viable on realistic edge hardware, rather than to
validate decision-making accuracy against a live water body. Every number in Section~\ref{sec:results}
comes from one of eleven experiments run directly against the S.U.R.E. codebase's own already-running
local configuration at git commit \texttt{3c1b9fa}, not from a synthetic benchmark constructed to
flatter the architecture. This section states, honestly, what each experiment measured, on what
sample, with what hardware, and what its own stated limitations are -- the same discipline the paper
asks of its own headline numbers is applied here to how they were obtained.

\subsection{General methodology and guardrails}

All experiments ran read-only against the running codebase and, where applicable, the live
\texttt{sure\_rag} database: no experiment inserted, updated, or deleted a database record, and
no experiment modified \texttt{sure-project} source files. Three experiments (RAG calibration,
vision leakage, and export benchmarking) required new, purpose-built analysis scripts rather than
reusing an existing one, because the existing scripts either wrote to a live database
(\texttt{rag/bench.py}) or the needed diagnostic did not exist yet (perceptual-hash leakage
detection, per-image ONNX-versus-TorchScript agreement). Every script and its raw output is
preserved alongside its experiment's \texttt{results.md} file for audit. Hardware for every vision
and edge-latency measurement was a single Apple M4 Pro; no claim in this paper generalizes across
device classes, and none is presented as if it did.

\subsection{Decision-fusion audit and behavioral trace (G2, G3)}

The claim that the evaluation harness and the system's decision path compute severity
identically was checked at the strongest available level: not by comparing outputs, but by
confirming, via Python's \texttt{is} identity operator, that both paths import and hold a reference
to the exact same in-memory \texttt{rules} module object. This was verified across all 8 of the
evaluation harness's built-in scenarios, with 8/8 agreement.

The dual-layer decision system's actual behavior was then traced across those same 8 scenarios
($n=8$, the harness's entire fixed scenario population -- this is not a subsample, and no
confidence interval is implied by or should be read into this count). For each scenario, the real,
unmodified \texttt{apply\_rule\_override} function was invoked against AQUA-1B's live output and
the rule engine's independently computed status, and each of the 8 outcomes was classified into one
of four mutually exclusive buckets (unparseable-defaulted-to-\texttt{ok}, parseable-and-agrees,
parseable-and-under-calls-and-escalated, parseable-and-over-calls). A separate mechanism-correctness
check (\texttt{backend/test\_decision.py}, 18/18 passing under synthetic fixtures) was run and is
reported strictly apart from this behavioral trace, per this project's own methodological
discipline: a unit test confirming the override logic is implemented correctly is not evidence
about how the real model's output actually distributes across that logic's branches, and the two
are never substituted for one another anywhere in this paper. All 8 free-text reasoning strings
were additionally hand-coded for factual accuracy against the true sensor snapshot each scenario
supplied -- a qualitative, single-rater judgment, reported as illustrative rather than a
statistically powered claim, but checkable in principle by any reader against the same 8 published
snapshots.

\subsection{Agentic tool-routing benchmark (G4)}

The agentic tool-calling benchmark ran two models (AQUA-1B and AQUA-7B, base weights, no LoRA
adapter loaded for either) against a closed menu of three read-only tools, first reproducing the
original 5-scenario benchmark, then extending it with 4 newly authored scenarios designed
specifically to stress a different failure mode than the original 5 already covered (a different
sensor parameter's trend, a recovering rather than degrading trend, a combined vision-and-sensor
cue, and a borderline-but-safe restraint case), for a total $n=9$ scenarios $\times$ 2 models = 18
independent agent runs. Greedy decoding (temperature 0) was used throughout, so no scenario was run
by simple repetition -- reported effect sizes are entirely a function of scenario diversity, not
resampling the same input.

\subsection{Vision dataset leakage audit (G5)}

Cross-split leakage between the 412-image training set and 98-image validation set was checked two
ways: an exact frame-index overlap check (trivial to compute, and confirming zero exact
duplicates), and a perceptual near-duplicate check using 64-bit average-hash and difference-hash
fingerprints computed for every image (the \texttt{imagehash} library was unavailable in this
environment; a minimal equivalent was implemented directly from PIL and NumPy). A known-positive
control -- a separate, already-documented dataset with training and validation splits pointing at
the identical folder -- was run through the same pipeline first, to confirm the method actually
detects a leak it is known to contain, before trusting a null or partial result on the real split.

\subsection{Export-format benchmark and edge latency (G6)}

Six export configurations of the same trained checkpoint (\texttt{best.pt}, epoch 77) were
benchmarked against the identical 98-image validation split: native PyTorch on both MPS and CPU,
ONNX (fp32 and INT8), CoreML (dispatched to the Apple Neural Engine), and TorchScript. Latency was
measured as p50/p95 over pre-decoded validation frames (avoiding disk I/O inside the timed loop),
discarding warm-up runs, following the same protocol the system's own export benchmarking script
uses by default. Because a single latency session can be affected by thermal state or scheduler
noise, the CoreML/ANE configuration -- the one actually configured for use -- was re-measured across three
independent process invocations rather than reported from a single run. A new per-image diff
between the ONNX and TorchScript exports' detections (greedy IoU-matched, 0.9 IoU bar for
agreement) was run across the full 98-image set to test directly whether their shared accuracy loss
reflects a shared mechanism rather than independent numerical noise.

\subsection{Vision operating-point safety trace (G7)}

The configured detection confidence threshold was confirmed directly from the vision service's own
source (\texttt{CONF\_THRESH=0.20}) rather than assumed from documentation. Two complementary
methods were used to characterize precision and recall as a function of confidence: sweeping
\texttt{model.val(conf=X)} over a coarse grid, and separately reading Ultralytics' underlying
fine-grained (1000-point) precision/recall/F1 curve directly from a single validation call. The
first method was found, over the course of this audit, to have a methodological trap worth
recording for other practitioners: it reports an argmax-over-the-retained-curve figure regardless
of the requested confidence, so it cannot answer ``what is precision/recall \emph{at} this specific
threshold'' -- only the second, curve-reading method can, and it is the one this paper's headline
numbers rely on. The full-frame-miss rate (a frame producing zero detections despite containing
ground-truth fish) was measured empirically across all 98 validation frames at the configured
threshold, and separately estimated via an independence approximation
$P(\text{miss}\mid k) = (1-\text{recall})^k$ for comparison. Because the validation set's minimum
ground-truth fish-per-frame count is $k=3$ -- it contains zero frames with $k=1$ or $k=2$ fish --
the sparse-frame regime most relevant to a welfare check involving a small or isolated fish could
only be estimated by extrapolation, never measured directly; this is reported as a dataset-coverage
gap, not folded silently into the headline result.

\subsection{PSI drift-binning sweep (G8)}

A 16-point severity sweep ($\delta \in \{0.00, 0.02, \ldots, 0.30\}$) was constructed by resampling
the reference confidence distribution ($n=1467$) with replacement and applying a mean shift of
$\delta$, clipped to $[0,1]$, then computing PSI under both quantile-edge and fixed equal-width
binning at every severity level. This is a mean-shift-only perturbation model; variance inflation
was not layered in and is noted as a natural extension rather than attempted here. The three-way
\texttt{decide()} gate and the promotion \texttt{gate()} function were each additionally exercised
directly against two hand-constructed synthetic windows (no-drift and significant-drift) and three
hand-constructed promotion-margin cases, to confirm the gate's documented behavior end-to-end.

\subsection{twin\_bridge mechanism-evidence inventory (G9)}

The existing PLC/CODESYS bridge test suite (\texttt{twin\_bridge/test\_bridge.py}) was run as-is,
and four new edge-case tests were authored against a scripted \texttt{FakeTwin} test double to
exercise a comparison-logic branch the existing suite never reached. No live Godot or CODESYS
session was available in this environment, and none was constructed for this audit; every test in
this subsystem, existing and new, is team-authored mechanism evidence run against a scripted
double, not evidence from an independently produced external system.

\subsection{RAG threshold calibration (G1) and manuscript-wide provenance audit (G10, G11)}

The RAG threshold sweep and the e5-small-versus-\texttt{tr-bert} retrieval comparison were
reproduced read-only against the live 44-chunk indexed corpus. Because the original benchmark
script clears and re-populates the live collection (a write operation forbidden under this audit's
read-only guardrail), a new script reproduced only the query-and-scoring half of that benchmark
against whatever was already indexed, after first confirming the indexed content matched the
current knowledge-base files exactly. A new random-retrieval baseline (uniform selection over all
44 chunks, evaluated across 5 seeds for stability) was authored specifically to establish that
e5-small's measured retrieval skill is a real effect and not an artifact of the corpus's small
size. Separately, a full manuscript-wide provenance and stale-figure audit re-verified every
headline number cited in this paper against its raw source file and git commit, and re-ran a
repository-wide grep for the specific superseded recall figure (0.695) this project's own history
required guarding against; that audit is what surfaced the still-unfixed occurrence of that figure
inside the system's RAG knowledge base reported in Section~\ref{sec:results}.
